\begin{proposition}\label{prop:all-octic-39}
Every rational effective divisor of degree $8$ on $X_1(39)$ supported
away from the cusps is the closed octic divisor represented by the
known CM point in the realization table. Consequently an elliptic
curve over an octic number field with a point of order $39$ has full
torsion $C_{39}$. In particular, $C_{78}$, $C_{117}$ and
$C_3\times C_{39}$ do not belong to $\Phi(8)$.
\end{proposition}
\begin{proof}
Put $X=X_1(39)$. By \cite[Theorem~3.1(2), Theorem~4.13 and
Table~2]{DEHMZ}, its full rational Jacobian is rational cuspidal with
invariants $[7,31122,3236688]$. In particular,
$J_X(\mathbb Q)[2^\infty]\simeq C_2\times C_{16}$.
The rational gonality of $X$ is $14$ by \cite[Table~1]{DvH}.

We first prove that torsion specialization at $2$ is injective. Let
$H=\langle\langle16\rangle\rangle$, of order $3$, and let
$\pi:X\to Y=X/H$. The quotient has genus $9$ and good reduction at
$2$. Indeed, the source has good reduction, and this is a tame
quotient. The induced cyclic cover $Y\to X_0(39)$ has degree $4$
and is etale. At an elliptic point the only possible nontrivial
stabilizer has order $3$, which is contained in $H$; order $4$ would
require a square root of $-1$ modulo $3$. Primitive cusp columns
show that the full diamond group acts freely on the cusp fibres.
The same argument applies in characteristic $2$.

Here are independent exact checks of the special-fibre model.
The standard model of $X_0(39)$ has genus $3$, a degree-$56$
$j$-map, and four rational cusps of widths $1,3,13,39$.
Its Picard group over $\mathbb F_2$ is $C_{28}\times\mathbb Z$.
There is a unique character to $C_4$ with Frobenius exponents
$(1,3,2,0)$ in that order, relative to the inverse of the usual
diamond generator. Inverting all exponents leaves the defining
kernel unchanged. The associated length-two Witt tower is converted
to an explicit quartic with exact inverse maps. Its deck map has
order $4$, and its square translates the quartic generator by $1$.
Witt traces on generators of the full Picard group modulo $4$
verify the character. The resulting curve has genus $9$ and a split
rational cusp; Riemann--Hurwitz verifies that the cover is etale.
Global class field theory therefore identifies it with
$Y_{\mathbb F_2}$. The primitive-column computation gives identical
rational and finite-field cusp-orbit degrees $1^4,2^2,4^2$.
In particular every closed cusp used in the computation lifts to a
rational cuspidal divisor.

The exact Picard group of this model is
$C_6\times C_{336}\times\mathbb Z$. The reductions of rational
cuspidal classes generate its full degree-zero subgroup. Thus their
$2$-primary image has order $32$. On the other hand, $\pi^*$ is
injective on rational $2$-primary torsion, since
$\pi_*\pi^*=[3]$. It follows that $J_Y(\mathbb Q)[2^\infty]$
has order at most $32$, and the computed image attains this bound.
Specialization is therefore injective on that subgroup.

Every automorphism of $C_2\times C_{16}$ sends generators $e,f$ to
$e+\epsilon8f$ and $\eta e+u f$, respectively, where
$\epsilon,\eta\in\{0,1\}$ and $u$ is odd modulo $16$.
Its automorphism group consequently has order $32$. Hence the
order-three group $H$ acts trivially on $J_X(\mathbb Q)[2^\infty]$.
The identity $\pi^*\pi_* = \sum_{h\in H}h=[3]$ makes $\pi_*$
injective there. Compatibility with specialization and the preceding
target injectivity prove source injectivity on rational $2$-primary
torsion. The odd-primary torsion is automatically injective at $2$.
This proves injectivity on all of $J_X(\mathbb Q)$.
The actual Picard relation matrix and cusp vectors are exported;
an independent integer Hermite normal form verifies the full cusp
image. A separate enumeration verifies the order-$32$ automorphism
group and all source and target cusp orbits.

Let $D$ be a rational noncuspidal effective divisor of degree $8$.
The source has the same rational and $\mathbb F_2$ cusp degrees
$1^{12},2^6,12^2$. The modular-unit signatures separate the rational
orbits, and primitive columns independently show that each remains
irreducible. These signatures commute with reduction: the functions
$b$ and $b^4/\Delta$ are units on the integral fine open modular
curve, with divisors supported on its etale cuspidal boundary.
Thus entirely cuspidal reduction lifts to a rational effective
cuspidal divisor $C$ of degree $8$. Injectivity would give $D\sim C$,
contradicting gonality $14>8$.

The complete noncuspidal place counts in degrees one through eight are
\[
 (0,0,0,0,0,12,0,49).
\]
They are independently checked by Newton identities from the
modular-symbol Frobenius polynomial. Hence a remaining reduction
consists of a sextic place with a cuspidal complement of degree $2$,
or of a single octic place. The exact Tate-coordinate diamond
partition gives one sextic orbit and five octic orbits.
All possible cuspidal complements lift to rational divisors.
Put $A_5=T_5-1-5\langle5\rangle$. Eichler--Shimura and injectivity
of torsion reduction at the odd prime $5$ give
$A_5J_X(\mathbb Q)=0$. On each rational cuspidal divisor $C$,
$T_5C=C+5\langle5\rangle C$, so $A_5C=0$ as a divisor.
The fresh computation proves that $A_5P$ is nonprincipal for the
sextic representative and for four octic representatives. The fifth
octic representative is exactly the reduction of the known point.
Its full finite diamond orbit is a singleton. Thus these five
nonprincipality witnesses cover every unknown noncuspidal reduction.

The known point has $j=0$ and is defined over the irreducible octic field
\[
 x^8-4x^7+10x^6-25x^5+37x^4-25x^3+10x^2-4x+1=0.
\]
Its universal Tate curve and second coordinate are given in the
realization table. This polynomial remains irreducible modulo $2$,
and its coordinate reduction is the surviving octic place $\overline E$.
Consequently $\overline D=\overline E$. Injectivity gives $D\sim E$,
and gonality $14>8$ gives $D=E$.

The known curve has exact point order $39$ and good unramified
reduction of order $39$ over $\mathbb F_{49}$, so its full torsion is
$C_{39}$. Finally, apply the divisor classification to the rational
degree-eight trace divisor of a moduli point over an octic field,
with multiplicities if its field of definition is smaller. The
divisor is forced to be closed octic. The fine moduli interpretation
then identifies the elliptic curve over that field with the known
model. This proves the asserted full-torsion classification.
\end{proof}
